% Evidence-backed replacement candidate, 2026-10-18.
% Intentionally not inserted into the legacy SLAST manuscript automatically:
% task, abstract, method, and figure must be revised consistently first.
\section{单风场预测与跨风场适配结果}
\label{sec:single-farm-review}

评价目标为单风场未来六小时、十五分钟分辨率的功率轨迹。
源站任务采用56站、2022年训练、2023年验证和2024年跨年复评。
训练可共享参数，部署时的本站输入与区域历史辅助分别报告。
主指标为每站容量因子MAE，先站内计算、再等权平均风场和随机种子；
均值是单模型误差的平均，不是集成模型误差。
各年份采用全部有效缓存窗口，并非全年无缺测。

\begin{table}[tbp]
\centering\small
\caption{源56站的单风场预测。十种子均值；Persistence无训练。
区域历史版本与本站输入版本是系统对照，不是同输入的单变量消融。}
\begin{tabular}{llrr}
\toprule
模型 & 在线输入 & 2023 MAE & 2024 MAE\\
\midrule
Persistence & 本站历史 & 0.117542 & 0.112763\\
共享Transformer & 本站历史/天气 & 0.087732 & 0.084699\\
CBH-SL & 本站历史/天气 & 0.084140 & 0.081799\\
CBH-IL（独立训练） & 本站历史/天气 & 0.088450 & 0.086998\\
\midrule
CBH-SC & 加区域历史 & 0.083814 & 0.081507\\
E0-CBH（FP32） & 加区域历史 & 0.083827 & 0.081528\\
E2C-CBH（FP32） & 加区域历史 & 0.082667 & 0.080703\\
\bottomrule
\end{tabular}
\end{table}

CBH-SL相对共享Transformer在2023/2024年的MAE降低4.09\%/3.42\%，
2024年在49/56站更好；场站配对bootstrap区间为
$[-0.003651,-0.002095]$。
SL相对IL的2024年差值为$-0.005199$，54/56站改善，
支持不依赖在线区域功率的共享学习。
E2C在区域历史设置下进一步改善，但不将其直接称为本站输入部署模型。
源站独立命名baseline仅使用月抽样pilot，训练/验证窗口为3156/3349，
因此该批结果另表报告，不进入全年主表排名。

\begin{table}[tbp]
\centering\small
\caption{重新训练的FP16条件消融。各臂保留公共历史记忆和输出基准。
CB仍包含$b$与$\Delta b$；CBH相对CB仅增加额外历史摘要。}
\begin{tabular}{lccc rr}
\toprule
变体 & 摘要 & $b$ & $\Delta b$ & 2023 MAE & 2024 MAE\\
\midrule
CBH & 有 & 有 & 有 & 0.083781 & 0.081471\\
CB & 无 & 有 & 有 & 0.083961 & 0.081731\\
CH & 有 & 无 & 无 & 0.088499 & 0.086065\\
CNEUTRAL & 无 & 无 & 无 & 0.088549 & 0.086151\\
\bottomrule
\end{tabular}
\end{table}

CB相对CNEUTRAL的2024年差值为$-0.004419$，支持轨迹条件信息包。
CBH相对CB的差值为$-0.000260$，表明摘要提供较小增量。
四臂未单独切换$b$与$\Delta b$，不能将两者完全分离归因。
匹配的direct/correction对照差约$0.000096$，区间跨零，
故纯残差输出锚点不足以解释完整结构收益。
天气编码器的等参数E2N/E2L控制支持非线性激活，
但DEM、柱坐标和显式分组未形成独立稳健优势。

区域历史版本的分时效结果显示，E2C相对E0在2024年0--1h、1--3h和3--6h
窗口分别改善约$0.000159$、$0.000620$和$0.001185$，春季改善$0.001355$。
未来真实一小时变化超过0.2 CF的窗口中改善$0.001892$；
该切片使用未来真值，只作误差诊断，不作为在线模型选择条件。
源站案例采用2022年运行描述量确定的六个代表站，并保留龙湖失败案例。

\begin{table}[tbp]
\centering\small
\caption{20个未见目标站的2024 H2评价，共同五种子。
预算仅计适配训练标签，额外使用完整2024 H1标签选checkpoint。
Scratch家族在dev10的H1上冻结。}
\begin{tabular}{llrrr}
\toprule
预算 & Scratch家族 & Scratch MAE & 预训练Transformer & E2C适配\\
\midrule
7d & GRU & 0.123650 & 0.094832 & 0.095585\\
30d & TFT & 0.124890 & 0.093827 & 0.095062\\
90d & TimeXer & 0.120178 & 0.092754 & 0.093435\\
180d & Transformer & 0.108321 & 0.092510 & 0.092528\\
full & TiDE & 0.096610 & 0.090381 & 0.089853\\
\bottomrule
\end{tabular}
\end{table}

E2C适配相对dev选定scratch的收益明显，但30d下相对预训练Transformer
的差值为$+0.001235$，场站95\%区间为$[+0.000049,+0.002509]$；
full的平均优势区间仍跨零。因此不能声称条件表示比通用Transformer更可迁移。
现有E2C迁移采用FP32且读取源站聚合历史，与历史对照的精度/输入并非完全匹配。
严格总标签预算和本站输入的迁移仍需补充实验。

统计上，原E2C报告以有效origin加权MAE为点估计，却对七日块等权构造区间。
2024年有效origin权重下差值$-0.000825$，seed、风场与共同七日块重采样区间为
$[-0.001358,-0.000292]$；等周敏感性估计为$-0.000557$，区间为
$[-0.001178,+0.000168]$。两者统计目标不同，应同时呈现。
2024年已参与多轮复评，不称为未触碰盲测。
